\documentclass[10pt,a4paper]{amsart}
\usepackage[margin=19mm]{geometry}
\usepackage{amsmath,amssymb,mathtools}
\usepackage[scr=rsfs,cal=euler]{mathalfa}
\usepackage{xfrac}
\usepackage[unicode,hidelinks,bookmarks=false]{hyperref}
\newcommand{\ie}{i.e.\ }
\newcommand{\cf}{cf.\ }
\newcommand{\resp}{respectively\ }
\newcommand{\Subsection}{\subsection}
\providecommand{\nobreakdash}{\nobreak\hbox{-}}
\setlength{\emergencystretch}{2em}
\multlinegap0pt
\usepackage{amssymb}
\usepackage{bm}
\newtheorem{lemma}{Lemma}
\title{On Permutation Groups of Cyclic Codes over Finite Fields}
\author{Junjie Huang, Jicheng Ma, Chang-An Zhao}
\date{Source: arXiv:2605.24314v1 (CC BY 4.0)}
\begin{document}
\maketitle

\noindent\emph{Source and licence.} On Permutation Groups of Cyclic Codes over Finite Fields. Version arXiv:2605.24314v1 (CC BY 4.0), distributed by arXiv under the Creative Commons Attribution 4.0 International licence, http://creativecommons.org/licenses/by/4.0/, as recorded in the arXiv licence metadata for this exact version and shown on the abstract page https://arxiv.org/abs/2605.24314v1. Full licence notice: \texttt{licenses/arxiv-2605.24314-CC-BY-4.0.txt}.\par\smallskip

\noindent\textit{Editorial scope.} Counted unit: the article's lemma L\_hp\_col (source0.tex lines 274--276). The article's complete original proof of it is delivered with it (source0.tex lines 277--376, one \texttt{proof} environment of 100 lines). No mathematics is written, repaired or re-derived: the statement and the proof are the article's own lines, in the article's own notation, and no retained line is substituted or re-flowed.  Retained uncounted as the source-local context the proof needs: lines 254--256.  Nothing was omitted from any retained span, so the package carries no omission record.  No retained line contains a citation, so the wrapper carries no bibliography and no bibliography entry.  No retained line contains a figure environment, an image-inclusion command or drawing code, so no image is delivered and the package declares no asset dependency.  Editorial formatting only, in addition: an AMS \texttt{amsart} preamble that restates the article's own declarations and macros used by the retained lines, namely the environments \texttt{lemma} on separate counters, together with the standard packages those lines need.  The retained environments print in this document as Lemma 1 in order of appearance, whereas the article prints the same items with the numbering of its own full document.  The complete native source of the article as distributed in the arXiv e-print for this exact version is bundled unchanged as \texttt{sources/201264/source0.tex}.
\begin{lemma}\label{L_hp_row}
    The permutation group $\operatorname{Per}(\mathcal{C}_{hp,g(x)})$ of $\mathcal{C}_{hp,g(x)}$ contains a subgroup that is isomorphic to $(S_h)^p$. 
\end{lemma}

\begin{lemma}\label{L_hp_col}
    The permutation group $\operatorname{Per}(\mathcal{C}_{hp,g(x)})$ of $\mathcal{C}_{hp,g(x)}$ contains a subgroup $G$, where $G\cong \operatorname{Per}(\mathcal{C}_{p,g(x)})$, the permutation group of $\mathcal{C}_{p,g(x)}$. 
\end{lemma}

\begin{proof}
    Let $c$ be any generator codeword of $\mathcal{C}_{hp,g(x)}$. For each permutation $\sigma\in S_p$, let $\bar{\sigma}\in S_{hp}$ be the lifting permutation of $\sigma$ such that $\bar{\sigma} = \rho_0(\sigma)\rho_p(\sigma)\rho_{2p}(\sigma)\cdots\rho_{(h-1)p}(\sigma)$ where $\rho_{ip}$ maps coordinates $j$ in $\sigma$ to $ip+j$ with $0\le i \le h-1$ and $1\le j \le p$. Indeed, the permutation $\bar{\sigma}$ is the column permutation of $Mr_{h}(c)$. Now, we claim that $\bar{\sigma}\in \operatorname{Per}(\mathcal{C}_{hp,g(x)})$ if and only if $\sigma\in \operatorname{Per}(\mathcal{C}_{p,g(x)})$. Firstly, it is easy to verify that if
    \[
        Mr_{h}(a^\prime) = \begin{pmatrix}
            \cdots \\
            {\bf 0}_p \\
            a \\
            {\bf 0}_p \\
            \cdots
        \end{pmatrix}
    \]
    for any $a^\prime\in\mathcal{C}_{hp,g(x)}$, then $a\in \mathcal{C}_{p,g(x)}$ by considering the linear combination of the basis of $\mathcal{C}_{hp,g(x)}$. On the contrary, if $a\in \mathcal{C}_{p,g(x)}$, then it is easy to see that
    \[
        \begin{pmatrix}
            \cdots \\
            {\bf 0}_p \\
            a \\
            {\bf 0}_p \\
            \cdots
        \end{pmatrix} = Mr_{h}((0\cdots,0,a,0,\cdots,0))\in\mathcal{C}_{hp,g(x)}.
    \]

    Assume that $\bar{\sigma}\in \operatorname{Per}(\mathcal{C}_{hp,g(x)})$. Then for any $a\in \mathcal{C}_{p,g(x)}$, we have
    \[
        Mr_{h}((a,0,\cdots,0))^{\bar{\sigma}} = \begin{pmatrix}
            a^\sigma \\
            {\bf 0}_p \\
            \cdots \\
            {\bf 0}_p
        \end{pmatrix}\in\mathcal{C}_{hp,g(x)}.
    \]
    Therefore, $a^\sigma\in\mathcal{C}_{p,g(x)}$. It follows that $\sigma\in \operatorname{Per}(\mathcal{C}_{p,g(x)})$. 

    Now, we assume that $\sigma\in \operatorname{Per}(\mathcal{C}_{p,g(x)})$. 
    \begin{itemize}
        \item If $Mr_{h}(c)$ is of the form
        \[
            \begin{pmatrix}
                \cdots \\
                {\bf 0}_p \\
                v_0^{\sigma_0^t} \\
                {\bf 0}_p \\
                \cdots
            \end{pmatrix},
        \]
        then we have
        \[
            Mr_{h}(c)^{\bar{\sigma}} = \begin{pmatrix}
                \cdots \\
                {\bf 0}_p \\
                v_0^{\sigma_0^t\cdot\sigma} \\
                {\bf 0}_p \\
                \cdots
            \end{pmatrix}.
        \]
        Note that $v_0^{\sigma_0^t\cdot\sigma}\in\mathcal{C}_{p,g(x)}$, therefore it can be linearly represented by $v_0,v_0^{\sigma_0},\cdots,v_0^{\sigma_0^{p-m-1}}$. Hence $Mr_{h}(c)^{\bar{\sigma}}\in \mathcal{C}_{hp,g(x)}$. 
        \item If $Mr_{h}(c)$ is of the form
        \[
            \begin{pmatrix}
                \cdots \\
                {\bf 0}_p \\
                v_1 \\
                v_2 \\
                {\bf 0}_p \\
                \cdots
            \end{pmatrix},
        \]
        then we have
        \[
            Mr_{h}(c)^{\bar{\sigma}} = \begin{pmatrix}
                \cdots \\
                {\bf 0}_p \\
                v_1^{\rho_{ip}(\sigma)} \\
                v_2^{\rho_{(i+1)p}(\sigma)} \\
                {\bf 0}_p \\
                \cdots
            \end{pmatrix},
        \]
        for some $0\le i \le h-1$. Since $v_1+v_2 = v_0^{\sigma_0^l}$ for some $p-m\le l\le p-1$, we have $v_1^{\rho_{ip}(\sigma)} + v_2^{\rho_{(i+1)p}(\sigma)} = v_0^{\sigma_0^l\cdot \sigma}\in \mathcal{C}_{p,g(x)}$. Moreover, we have 
        \[
            Mr_{h}(c)^{\bar{\sigma}} - \begin{pmatrix}
                \cdots \\
                {\bf 0}_p \\
                v_1^{\rho_{ip}(\sigma)}+v_2^{\rho_{(i+1)p}(\sigma)} \\
                {\bf 0}_p \\
                {\bf 0}_p \\
                \cdots
            \end{pmatrix} = \begin{pmatrix}
                \cdots \\
                {\bf 0}_p \\
                -v_2^{\rho_{(i+1)p}(\sigma)} \\
                v_2^{\rho_{(i+1)p}(\sigma)} \\
                {\bf 0}_p \\
                \cdots
            \end{pmatrix}\in \mathcal{C}_{hp,g(x)}\,(\text{Similar to the proof of Lemma \ref{L_hp_row}}).
        \]
        Hence, we derive that $Mr_{h}(c)^{\bar{\sigma}}\in \mathcal{C}_{hp,g(x)}$.
    \end{itemize}
    From the above deduction, we conclude that $\bar{\sigma}\in \operatorname{Per}(\mathcal{C}_{hp,g(x)})$. It follows that all permutation $\bar{\sigma}\in\operatorname{Per}(\mathcal{C}_{hp,g(x)})$ that permute columns generate a subgroup $G$ which is isomorphic to $\operatorname{Per}(\mathcal{C}_{p,g(x)})$. The proof is completed.
\end{proof}

\end{document}
